\documentclass[10pt,a4paper]{amsart}
\usepackage[margin=19mm]{geometry}
\usepackage{amsmath,amssymb,mathtools}
\usepackage[scr=rsfs,cal=euler]{mathalfa}
\usepackage{xfrac}
\usepackage[unicode,hidelinks,bookmarks=false]{hyperref}
\newcommand{\ie}{i.e.\ }
\newcommand{\cf}{cf.\ }
\newcommand{\resp}{respectively\ }
\newcommand{\Subsection}{\subsection}
\providecommand{\nobreakdash}{\nobreak\hbox{-}}
\setlength{\emergencystretch}{2em}
\multlinegap0pt
\usepackage{bm}
\usepackage{bbm}
\usepackage{dsfont}
\usepackage{xspace}
\usepackage{cancel}
\usepackage{mathtools}
\usepackage{amsthm}
\makeatletter
\@ifundefined{theorem}{\newtheorem{theorem}{Theorem}}{}
\makeatother
\title[A Complete Classification of Low-Order Conservation Laws for]{A Complete Classification of Low-Order Conservation Laws for a Generalized Fifth-Order KP Family}
\author{Nitin Serwa}
\date{Source: arXiv:2605.02473v2 (CC BY 4.0)}
\begin{document}
\maketitle

\noindent\emph{Source and licence.} A Complete Classification of Low-Order Conservation Laws for a Generalized Fifth-Order KP Family. Version arXiv:2605.02473v2 (CC BY 4.0), distributed under the Creative Commons Attribution 4.0 International licence, as recorded in the arXiv licence metadata for this exact version and shown on the abstract page https://arxiv.org/abs/2605.02473v2. Full rights notice: licenses/arxiv-2605.02473-CC-BY-4.0.txt.\par\smallskip

\noindent\textit{Editorial scope.} Counted unit: the Theorem whose source label is \texttt{thm:order-reduction} in the article (source0.tex lines 624--640, source label \texttt{thm:order-reduction}), delivered together with the article's own complete original proof of it (source0.tex lines 642--682, one \texttt{proof} environment of 41 lines). No mathematics is written, repaired or re-derived: statement and proof are the article's own text line for line. Required context retained uncounted: eq:normalized-equation (lines 376--383). Every definition, display and label of those lines is retained unchanged. Omitted from this package and not redistributed: 1--375, 384--623, 641--641, 683--1784. Nothing inside a retained excerpt is omitted or altered: every retained body is delivered exactly as the article's lines are written, so no omission receipt of any kind accompanies this package. Citations: the retained excerpts carry no citation command at all, so no bibliography entry of the article is transcribed and no reference is redistributed. Reference adaptation: none was needed and none was made; every reference inside the delivered unit resolves inside it, so no unresolved reference or citation is produced. Printed slips kept literal: no printed word, symbol, hypothesis or number of the counted statement or of its proof was altered. Layout and class: the article's own class is \texttt{article} with packages including inputenc, fontenc, lmodern, microtype, amsmath, amssymb, amsthm, mathtools, geometry, enumitem, hyperref; the wrapper is the lane's pinned \texttt{amsart} editorial wrapper at 10pt on A4, which restates the theorem-like environments the retained text needs and adds the article's own macro definitions that the retained text needs (none), with extra wrapper packages (bm, bbm, dsfont, xspace, cancel, mathtools). The delivered numbering is the unit's own. Assets: no retained span contains a figure environment, a graphics-inclusion command, a PDF-page inclusion, a table or a TikZ picture; no image file is copied into this package, the package declares no asset dependency and no image file is redistributed with it. Licence: version arXiv:2605.02473v2 of this article is distributed under the Creative Commons Attribution 4.0 International licence, as recorded in the arXiv licence metadata for this exact version and shown on the abstract page https://arxiv.org/abs/2605.02473v2. A full rights notice is delivered at licenses/arxiv-2605.02473-CC-BY-4.0.txt. Characters: every retained excerpt is pure ASCII, so the delivered engine, XeLaTeX with its default Latin Modern text font, reports no missing character.
\begin{equation}
\bigl(u_t+c\,u_{xxx}+u_{xxxxx}
+e\,u\,u_{xxx}+f\,u_xu_{xx}+g\,u^2u_x\bigr)_x
+\sigma u_{yy}=0,
\qquad
\sigma=\pm1.
\label{eq:normalized-equation}
\end{equation}

\begin{theorem}
\label{thm:order-reduction}
Let $Q$ be a multiplier of \eqref{eq:normalized-equation} of the form
\begin{equation}
Q = Q\bigl(x,y,t,u,u_x,u_y,u_t,
u_{xx},u_{xy},u_{yy},u_{yt},u_{tt}\bigr).
\label{eq:second-order-ansatz}
\end{equation}
Then $Q$ is independent of every second-order jet:
\begin{equation}
Q_{u_{xx}} = Q_{u_{xy}} = Q_{u_{yy}}
= Q_{u_{yt}} = Q_{u_{tt}} = 0,
\label{eq:no-second-order}
\end{equation}
so that $Q = Q(x,y,t,u,u_x,u_y,u_t)$. This holds for all
$c,e,f,g$ and $\sigma=\pm1$.
\end{theorem}

\begin{proof}
The determining equation $E_u(Q\Delta)=0$ is an identity in the jet
variables. Let $u_J$ be any one of the five second-order jets in
\eqref{eq:second-order-ansatz}, so that $|J|=2$, and consider in
$E_u(Q\Delta)$ the coefficient of the single jet $u_{J+(6,0,0)}$,
which has total order eight.

Two contributions to this coefficient arise. The Euler operator
contains the term
\[
(-D)_J\left[
\frac{\partial(Q\Delta)}{\partial u_J}
\right].
\]
Since $u_{xxxxxx}$ occurs in $\Delta$ linearly with coefficient one,
$\partial(Q\Delta)/\partial u_J$ contains $Q_{u_J}\,u_{xxxxxx}$, and
applying $(-D)_J$ produces $Q_{u_J}\,u_{J+(6,0,0)}$ with sign
$(-1)^{|J|}=+1$. The Euler operator also contains
\[
(-D_x)^6\left[
\frac{\partial(Q\Delta)}{\partial u_{xxxxxx}}
\right]
=
(-D_x)^6[Q].
\] The term in which all six derivatives act on $Q$
through its dependence on $u_J$ produces $Q_{u_J}\,u_{J+(6,0,0)}$ with
sign $(-1)^6=+1$.

No further contribution to the single jet $u_{J+(6,0,0)}$ occurs.
Since $Q$ has differential order at most two, it can supply a jet of
order eight only through a single differentiation of its dependence on
a second-order jet, as in the two contributions above; any other
Euler term either acts on a jet of $\Delta$ of $x$-order at most four,
yielding single jets of total order at most six, or arises from a
quadratic term of $\Delta$, yielding products of jets rather than the
single jet $u_{J+(6,0,0)}$. Hence the coefficient of $u_{J+(6,0,0)}$
equals $2Q_{u_J}$, and the determining equation forces $Q_{u_J}=0$.
Applying this to each of the five second-order jets gives
\eqref{eq:no-second-order}. No division by the coefficients is used,
so the conclusion is uniform in $c,e,f,g,\sigma$.
\end{proof}

\end{document}
